\documentclass[tikz,border=8pt]{standalone}
\usepackage{ctex}
\usepackage{tikz}
\usetikzlibrary{arrows.meta,positioning,fit,backgrounds}

\begin{document}
\begin{tikzpicture}[
  font=\small,
  >=Latex,
  blk/.style={draw=#1!70!black, fill=#1!8, rounded corners=2.2mm, very thick, minimum width=34mm, minimum height=10mm, align=center},
  flow/.style={->, very thick, draw=black!75},
  badge/.style={draw=black!35, rounded corners=1.5mm, fill=yellow!14, inner sep=3pt}
]

\node[blk=blue] (daemon) at (0,0) {real\_ids\_daemon\\64 帧历史窗口};
\node[blk=teal] (prep) at (4.8,0) {在线预处理\\64x9 + \(\Delta t\) 归一化};
\node[blk=purple] (infer) at (9.6,0) {CAN-CNN 前向推理\\\texttt{best\_model.pth}};
\node[blk=green!70!black] (fuse) at (14.4,0) {融合输出\\\texttt{can\_ml}};

\draw[flow] (daemon) -- (prep);
\draw[flow] (prep) -- (infer);
\draw[flow] (infer) -- (fuse);

\node[badge] at (4.8,-1.6) {\texttt{fixed preprocessing rule}};
\node[badge] at (9.6,-1.6) {\texttt{lightweight inference}};
\node[badge] at (14.4,-1.6) {\texttt{deployment ready}};

\node[draw=black!35, fill=orange!10, rounded corners=2mm, align=left, inner sep=5pt] at (7.2,2.0) {
\textbf{效率要点：}\\
1) 主路径仅窗口构造 + 单次前向；\\
2) 训练重计算离线完成；\\
3) 运行期内存与时延可控。};

\begin{scope}[on background layer]
  \node[draw=black!25, rounded corners=3mm, line width=0.8pt, fit=(daemon)(prep)(infer)(fuse), inner sep=7mm,
  label={[yshift=2mm]\large\bfseries 在线部署与效率（重绘版）}] {};
\end{scope}
\end{tikzpicture}
\end{document}
